\chapter*{Preface}
\addcontentsline{toc}{chapter}{Preface}
This book is a self-contained account of the Lean~4 library \textsc{LeanVOA} (this release edition
omits one chapter, on affine $\mathfrak{sl}_3$, which is in preparation): a
formalization of the theory of vertex operator algebras that culminates in a machine-checked
construction of the Moonshine module $\Vnat$ of Frenkel, Lepowsky and Meurman. It is written for a
reader who is an advanced graduate student in mathematics or mathematical physics, who knows
algebra and some representation theory, but who is neither an expert in vertex algebras nor a
user of Lean.

\paragraph{What the book is for.} A formal proof is only as good as its statement. A Lean theorem
can be perfectly verified by the computer and still fail to say what its authors meant: a
hypothesis may be too strong, a definition may be vacuous, a convention may have silently changed.
The purpose of this book is to let a \emph{human referee} check, theorem by theorem, that the Lean
statements are the mathematics we intended. For every definition and theorem we therefore print
(i) the statement in ordinary mathematical language, (ii) the Lean declaration copied verbatim
from the source, (iii) a \emph{fidelity check} that decodes the Lean statement and compares it with
the mathematics, (iv) a proof explained on paper, and (v) the machine status of the result. Each
chapter ends with a \emph{referee's checklist}. We do not explain how the Lean proofs are
organised: the computer has checked them; the reader's job, and ours, is to check that what was
checked is what we wanted.

\paragraph{How the book is tied to the code.} This edition documents \bookrelease{} of the library (commit
\texttt{\bookcommit}). The chapters do not repeat this: whenever a chapter speaks of ``the release'' or ``the released
source'', it means this commit, and every check that was run at an earlier state of the library and not repeated for this
edition is marked as such where it is reported. Every Lean file of the library is assigned to exactly one chapter of the full book (in this
edition, files assigned to the chapters in preparation are not documented), and every quotation of Lean
code in the book matches the source verbatim.

\paragraph{The trust base.} Lean checks every proof against its kernel. The library uses only Lean's three
standard axioms (\lean{propext}, \lean{Classical.choice}, \lean{Quot.sound}); it contains no project axiom.
Until 19 October 2026 it contained one, \lean{VOA.pbw}, a
Poincaré--Birkhoff--Witt statement for affine vacuum modules and induced modules; that declaration is now
a theorem with the same name and statement (Chapter~\ref{ch:L02}). Some chapters still record, as history,
which results used to depend on it; the construction of the Moonshine module never did. A clean axiom list
does not make a theorem unconditional: \texttt{\#print axioms} does not report the explicit hypotheses of a
theorem, so every status box of the book names them, and several headline results of Part~VI are
conditional reductions to named, unproved inputs (Chapter~\ref{ch:01}, Section~\ref{ch01:sec:newstatus}).

\paragraph{Defects found while writing the book.} Writing and refereeing the book was itself an audit
of the library. Every defect it found in the Lean source (a vacuous hypothesis, an uninhabitable
structure, a missing axiom field, a stale docstring) is reported in the chapter where it occurs,
in a box headed ``Trap / dead end'', and collected, with its evidence, in the catalogue of dead ends,
traps and documentation debt of Chapter~\ref{ch:15} (Sections~\ref{ch15:sec:deadends} and~\ref{ch15:sec:debt}). The book does not repair the library: it documents the release above
as it is.

\paragraph{Index.} The appendix ends with an index of every Lean declaration quoted here.

\paragraph{Structure.} Part~I sets up formal calculus, the axioms of vertex algebras and their
structure theory. Part~II treats the main families of examples and the rigidity programmes.
Parts~III and~IV develop representation theory and modularity. Part~V covers super, logarithmic and
unitary theories. Part~VI formalizes recent research: the Arakawa--Dai theorems on classical freeness,
Huang's tensor categories, the Gannon--Rayhaun classification of central charge one, affine vertex
algebras for every simple Lie algebra, and classical freeness for the Deligne exceptional series
(Chapters~\ref{ch:L24}--\ref{ch:L27}): a proof for $A_2$, and for the other six types a conditional
reduction whose remaining inputs are named; and the transfer of graded unitarity through the BRST coset
(Chapters~\ref{ch:L28} and~\ref{ch:L29}), from the boundary admissible affine $\mathfrak{sl}_2$ to the minimal models
$M(2,2n+1)$, a proved implication whose one hypothesis, graded unitarity of the affine algebra, is open. A chapter on
affine $\mathfrak{sl}_3$ is in preparation. Part~VII constructs the Moonshine module. Part~VIII records how all of this was
verified, every dead end met along the way, the numerical evidence, and the open problems. Chapter~\ref{ch:01},
Section~\ref{ch01:sec:map}, gives a map of the book, chapter by chapter, with the status of each headline result.
